\documentclass[11pt]{article}
\usepackage[autostyle, english = british]{csquotes}

% Small margins
\usepackage[top=0.75in, bottom=0.75in, left=0.75in, right=0.75in]{geometry}

% Better typography
\usepackage{mathptmx}  % Times-like font (professional and space-efficient)
\usepackage[scaled=0.92]{helvet}  % Sans-serif for headings
\usepackage{courier}  % Monospace for URLs/code if needed

% Line spacing for readability but compact
\usepackage{setspace}
\setstretch{1.1}

% No paragraph indentation, small skip instead
\setlength{\parindent}{0pt}
\setlength{\parskip}{6pt plus 2pt minus 1pt}

% Headers/footers with reviewer info
\usepackage{fancyhdr}
\pagestyle{fancy}
\fancyhf{}
\fancyhead[L]{\textit{Paper Review}}
\fancyhead[R]{\thepage}
\renewcommand{\headrulewidth}{0.4pt}

% Useful packages
\usepackage{url}
\usepackage{graphicx}
\usepackage{xcolor}

% For review comments
\usepackage[colorlinks, citecolor=blue, linkcolor=blue, urlcolor=blue]{hyperref}

% Simple review comment command
\newcommand{\reviewcomment}[2][red]{{\color{#1}\small [Reviewer: #2]}}

\begin{document}

% === Review header (compact) ===
\begin{center}
    {\large \textbf{NODLINK: An Online System for Fine-Grained APT Attack
Detection and Investigation}} \\[2pt]
    \textit{NDSS 2024} \quad | \quad \textit{Reviewer: Muhammad A'qil} \quad | \quad \textit{Date: \today}
\end{center}
\vspace{4pt}
\hrule
\vspace{8pt}

% === Introduction section (shortened) ===
\section*{Introduction}
NODLINK addresses real-time APT detection using provenance graphs. Existing online systems sacrifice granularity for speed, making investigation like ``searching for a needle in a haystack.'' The authors model detection as a \href{https://www.geeksforgeeks.org/dsa/steiner-tree/}{Steiner Tree Problem} (STP) and propose the \textbf{Importance-Score-Guided (ISG) greedy algorithm}, reducing complexity from $O(N^2)$ to $O(\theta N)$ by exploring top-$\theta$ important nodes per terminal.

Key features:
\begin{itemize}
    \item \textbf{In-memory cache:} prioritises anomalous nodes for long-running attacks
    \item \textbf{Terminal identification:} VAE-based anomaly detection
    \item \textbf{Hopset construction:} ISG for fast Steiner tree approximation
    \item \textbf{Comprehensive detection:} Grubb's test to flag attack campaigns
\end{itemize}

NODLINK was open-world evaluated across hospitals, universities, and factories, detecting 7 real attacks that HOLMES missed. Node-level precision is 2-3 orders of magnitude higher than UNICORN, with comparable throughput.

Contributions: (1) STP formulation for APT detection, (2) ISG algorithm with $O(\theta N)$ complexity, (3) commercial deployment, (4) first open-world evaluation for provenance-based APT detection.

% === Strengths section ===
\section*{Strengths}
\begin{itemize}
    \item \textbf{1. Novelty of Theoretical Framing:}
    NODLINK is the first work to formalise online APT detection as a STP with provable approximation bounds. This claim is supported by subsequent work (e.g. \href{https://www.sciencedirect.com/science/article/abs/pii/S1389128625006279}{CCLog}), which builds directly upon NODLINK's STP formulation and also cites the paper. The formal framing distinguishes NODLINK from prior heuristic-based systems (HOLMES, UNICORN).
    \item \textbf{2. Strong Evaluation Methodology:}
    First open-world evaluation of provenance-based APT detection systems, deployed in a live commerical SOC across 300+ machines across a variety of industries (hospitals, universities, and factories). NODLINK detected 7 real attacks that HOLMES missed entirely, while UNICORN produced 1000$\times$ more node-level false positives. Throughput is comparable to or better than SOTA systems, processing 40+ events per second per host.
    \item \textbf{3. Relevance and Impact:}
    NODLINK addresses a critical and urgent industrial problem -- APT detection under timeliness and resource constraints. It provides real benefits to security analysts, reducing the problem of "finding a needle in a haystack" significantly. The open source release and public datasets also lowers barriers for future research in similar provenance-based systems security. 
\end{itemize}

% === Weaknesses section ===
\section*{Weaknesses}
\begin{itemize}
    \item \textbf{1. ISG Algorithm Concerns} 
    \begin{itemize}
        \item \textbf{Formulation-implementation mismatch:} On page 5, the paper states the objective is to find an edge set $S$ with \textit{minimal total weight} connecting all terminals. However, the ISG algorithm (as defined in Section V-C) performs no minimisation in Hopset construction. It greedily explores nodes by importance score $IV(n)$ and stops after $\theta=10$ nodes per terminal. There is no step that compares alternative paths or selects a minimal-weight subset. The paper never reconciles this discrepancy.
        \item \textbf{Disconnected terminals not addressed:} The ISG algorithm terminates exploration after $\theta = 10$ hops per terminal. It exploits the "attack polymerism" property of APT attacks. If terminals are farther apart, they remain disconnected. The paper does not specify: (a) whether NODLINK raises alerts for disconnected hopsets, (b) whether it attempts to connect them across time windows, or (c) how often real attacks exceed this bound. Long-running APTs with extensive lateral movement across many hosts could potentially exceed this bound, rendering NODLINK unable to connect the full attack chain.
    \end{itemize}
    \item \textbf{2. Open-world ground truth limitations:} Ground truth relied on Sangfor's existing EDR rules to filter 135,700 alerts to 2,000 for manual review. This introduces potential bias toward attack patterns already known to Sangfor's rule set. Attacks not matching those rules would be absent from the ground truth. The paper does not comment on or attempt to estimate false negatives in the ground truth itself. More constructive independent open-world evaluation would involve third-party validation or red-team exercises unknown to the system.
    \item \textbf{3. No analysis of false negatives:} While the paper reports precision and recall, it does not analyze the specific attacks or attack steps that NODLINK misses. The authors note that "reconnaissance steps such as whoami, ipconfig, tasklist" are missed, but do not quantify how often this occurs or discuss whether such misses matter for attack campaign reconstruction.
\end{itemize}

% === Possible enhancements section ===
\section*{Possible Enhancements}
\begin{itemize}
    \item \textbf{$\theta=10$ adjustments:} Replace fixed $\theta=10$ with an adaptive mechanism. Some ideas I had: (a) gradually increasing $\theta$ until terminals connect till timeout is reached, (b) using graph diameter statistics to set $\theta$ per deployment environment, or (c) logging disconnected hopsets for analyst review. Specify fallback behavior when $\theta$ is insufficient.
    \item \textbf{Adversarial evaluation:} Evaluate NODLINK under targeted evasion strategies including: (a) gradient-based attacks against VAE embeddings, (b) high-IV decoy nodes designed to exhaust the $\theta$ budget, and (c) cache poisoning attacks. Report robustness metrics and discuss potential countermeasures.
\end{itemize}

% === Summary section ===
\section*{Summary}
NODLINK delivers a practically effective and open-world validated system for online APT detection. The system significantly reduces false positives compared to other SOTA systems, and the STP-inspired framing is a useful conceptual contribution.

However, several operational details (disconnected terminal behavior and ground truth construction) are underspecified. The authors should revise the theoretical framing to match the implementation and clarify the identified irregularities.

\end{document}